\chapter{Robust Statistics and Heavy Tails}\label{ch:qm:robust-statistics-and-heavy-tails}

One bad print in the feed, the euro's reference rate recorded as 1.5172 dollars instead of 1.1572 on one day in March 2026, multiplies
the standard deviation of the last 250 daily EUR/USD returns by seven, from 0.34\% to 2.43\%, and moves their \omterm{def:qm:robust-statistics-and-heavy-tails:huber}{median absolute
deviation} by 1\%. A risk system that estimates volatility with the standard deviation now reports a currency as volatile as a small
biotechnology stock; one that uses the \omterm{def:qm:robust-statistics-and-heavy-tails:huber}{median absolute deviation} has not noticed. The same data hold a second lesson: over 27 years
EUR/USD moved more than four standard deviations in a day seventeen times, where a normal law expects fewer than one. This chapter
is about \omterm{def:qm:estimation:estimator}{estimators} that bad data cannot capture and about the tails that good data do contain: the \omterm{def:qm:robust-statistics-and-heavy-tails:influence}{influence function} and the
\omterm{def:qm:robust-statistics-and-heavy-tails:influence}{breakdown point}, winsorising and ranks, \omterm{def:qm:robust-statistics-and-heavy-tails:copula}{copulas} and tail dependence, heavy tails and the \omterm{def:qm:robust-statistics-and-heavy-tails:gev}{extreme-value theory} that extrapolates them,
and the estimation of the \omterm{def:qm:robust-statistics-and-heavy-tails:heavy}{tail index} on real exchange rates.

\section{Robust losses and the influence function}

\begin{definition}[Influence function, breakdown point]\label{def:qm:robust-statistics-and-heavy-tails:influence}
For an \omterm{def:qm:estimation:estimator}{estimator} written as a functional $T(F)$ of the data's law, the \emph{influence function}\index{influence function} at $x$ is
$\mathrm{IF}(x) = \lim_{\varepsilon \to 0}\bigl(T((1 - \varepsilon)F + \varepsilon\delta_x) - T(F)\bigr)/\varepsilon$: the effect on the
estimate of an infinitesimal contamination at the point $x$ (Hampel, 1974). The \emph{breakdown point}\index{breakdown point} is the
smallest fraction of the sample that, replaced by arbitrary values, can carry the estimate arbitrarily far.
\end{definition}

The mean has $\mathrm{IF}(x) = x - \mu$, unbounded; the variance has $(x - \mu)^2 - \sigma^2$, worse. One observation placed far enough
moves either wherever it likes.

\begin{proposition}[Breakdown points of the mean and the median]\label{prop:qm:robust-statistics-and-heavy-tails:breakdown}
The sample mean has \omterm{def:qm:robust-statistics-and-heavy-tails:influence}{breakdown point} $1/n$, which tends to 0; the sample median has \omterm{def:qm:robust-statistics-and-heavy-tails:influence}{breakdown point} $\lfloor(n + 1)/2\rfloor/n$, which
tends to $1/2$, the largest possible for an equivariant location \omterm{def:qm:estimation:estimator}{estimator}.
\end{proposition}

\begin{proof}
Replacing one observation by $y$ moves the mean by $(y - x_1)/n$, unbounded in $y$. The median stays between the smallest and largest of
the uncontaminated observations as long as they are a majority; once half the sample is replaced, the contaminated values can be placed
on both sides of the median position and carry it. No equivariant \omterm{def:qm:estimation:estimator}{estimator} can do better: with half the sample shifted by $t$, the data
are equally consistent with a sample and with its shift, and the estimate cannot track both.
\end{proof}

\begin{definition}[Huber loss, median absolute deviation]\label{def:qm:robust-statistics-and-heavy-tails:huber}
The \emph{Huber loss}\index{Huber loss} with threshold $c$ is $\rho_c(r) = r^2/2$ for $|r| \le c$ and $c|r| - c^2/2$ otherwise; the Huber
\omterm{def:qm:estimation:m}{M-estimator} of location minimises $\sum_i\rho_c((x_i - \mu)/s)$ for a robust scale $s$, so its estimating function $\psi_c(r) = \max(-c,
\min(c, r))$ is bounded (Huber, 1964). The \emph{median absolute deviation}\index{median absolute deviation} is $\mathrm{MAD} =
1.4826\,\mathrm{med}_i|x_i - \mathrm{med}_jx_j|$, the constant making it consistent for the standard deviation of a normal law.
\end{definition}

The Huber \omterm{def:qm:estimation:estimator}{estimator} is the mean near the centre and the median in the tails. With $c = 1.345$ it keeps 95\% of the mean's efficiency when
the data are normal ($(2\Phi(c) - 1)^2/\E[\psi_c(Z)^2]$), and its \omterm{def:qm:robust-statistics-and-heavy-tails:influence}{influence function} is $\psi_c$, bounded by $c$ standard units. The MAD has
\omterm{def:qm:robust-statistics-and-heavy-tails:influence}{breakdown point} one half but is inefficient and assumes symmetry; the $Q_n$ \omterm{def:qm:estimation:estimator}{estimator} of Rousseeuw and Croux (1993), a scaled lower
quartile of all pairwise distances $|x_i - x_j|$, keeps the \omterm{def:qm:robust-statistics-and-heavy-tails:influence}{breakdown point} without the symmetry. \Cref{fig:qm:robust-statistics-and-heavy-tails:sensitivity}
shows what the \omterm{def:qm:robust-statistics-and-heavy-tails:influence}{influence functions} mean for one bad print: as the error on one day's log price grows from 0 to 27\%, the 250-day
standard deviation grows without limit, while the MAD and $Q_n$ move once, by 1.3\% and 1.5\%, and then not at all.

\begin{omfigure}
\begin{tikzpicture}
\begin{axis}[width=10cm,height=5cm,xlabel={error in one day's log price (\%)},ylabel={daily scale estimate (\%)},
  tick label style={font=\small},label style={font=\small},xmin=0,xmax=30,ymin=0,ymax=2.6,grid=major,grid style={gray!25},
  legend style={font=\footnotesize,at={(0.03,0.97)},anchor=north west,fill=white,draw=none},table/col sep=comma]
\addplot[omThm,very thick] table[x=error,y=sd]{figdata/methods/15-robust-statistics-and-heavy-tails/sensitivity.csv};
\addplot[omDef,very thick] table[x=error,y=mad]{figdata/methods/15-robust-statistics-and-heavy-tails/sensitivity.csv};
\addplot[black,thick,dashed] table[x=error,y=qn]{figdata/methods/15-robust-statistics-and-heavy-tails/sensitivity.csv};
\addplot[only marks,mark=*,mark size=1.8pt,black] coordinates {(27.09,2.4345)};
\legend{standard deviation,MAD,$Q_n$}
\end{axis}
\end{tikzpicture}

\omcaption{Three estimates of the daily volatility of EUR/USD from the 250 reference-rate returns to 23 September 2026, when the log
price of 24 March 2026 carries an error. The dot is the transposed-digit print 1.5172 for 1.1572 (an error of 27.1\%): standard
deviation 2.43\% instead of 0.34\%, MAD 0.279\% instead of 0.275\%, $Q_n$ 0.305\% instead of 0.300\%. Data: ECB euro reference rates
(source: ECB statistics).}\label{fig:qm:robust-statistics-and-heavy-tails:sensitivity}
\end{omfigure}

\section{Winsorising, trimming and rank methods}

\begin{definition}[Winsorisation, trimmed mean]\label{def:qm:robust-statistics-and-heavy-tails:winsor}
\emph{Winsorisation}\index{winsorisation} at level $p$ replaces the observations below the $p$-quantile by that quantile and those above the
$(1 - p)$-quantile by that one. The $p$-\emph{trimmed mean}\index{trimmed mean} discards the $\lfloor pn\rfloor$ smallest and largest
observations and averages the rest.
\end{definition}

Both have \omterm{def:qm:robust-statistics-and-heavy-tails:influence}{breakdown point} $p$. Winsorising is the common repair for a cross-section of signals or returns before a regression
(chapter~16): it keeps every observation's sign and rank and caps its leverage. It is also a decision about the data that the report must
state, because a winsorised \omterm{def:qm:estimation:sharpe}{Sharpe ratio} is not the strategy's \omterm{def:qm:estimation:sharpe}{Sharpe ratio}. For a bad print, the better repair is upstream: filtering
clearly erroneous trades against a robust band around a robust centre before any statistic is computed.

\begin{definition}[Rank correlation, Spearman's rho, Kendall's tau]\label{def:qm:robust-statistics-and-heavy-tails:rank}
A \emph{rank correlation}\index{rank correlation} measures dependence through the ranks of the observations only, so that it is unchanged
by increasing transformations of either variable. \emph{Spearman's rank correlation}\index{Spearman's rank correlation} is the correlation
of the ranks. \emph{Kendall's tau}\index{Kendall's tau} is the probability that two independent pairs are concordant minus the
probability that they are discordant, $\tau = \E[\operatorname{sign}((X - X')(Y - Y'))]$.
\end{definition}

\begin{definition}[Copula, tail dependence coefficient]\label{def:qm:robust-statistics-and-heavy-tails:copula}
A \emph{copula}\index{copula} is a joint distribution function on $[0, 1]^d$ with uniform margins. The lower \emph{tail dependence
coefficient}\index{tail dependence coefficient} of $(X, Y)$ with continuous margins $F_X, F_Y$ is $\lambda_L = \lim_{q \downarrow
0}\P(F_Y(Y) \le q \mid F_X(X) \le q)$, and the upper one is defined symmetrically.
\end{definition}

\begin{theorem}[Sklar]\label{thm:qm:robust-statistics-and-heavy-tails:sklar}
Every joint distribution function $H$ with margins $F_1, \dots, F_d$ can be written $H(x_1, \dots, x_d) = C(F_1(x_1), \dots, F_d(x_d))$ for a
\omterm{def:qm:robust-statistics-and-heavy-tails:copula}{copula} $C$, unique when the margins are continuous; conversely, any \omterm{def:qm:robust-statistics-and-heavy-tails:copula}{copula} and any margins define a joint law this way.
\end{theorem}

For continuous margins the proof is one line: $C$ is the law of $(F_1(X_1), \dots, F_d(X_d))$, whose margins are uniform. The theorem
separates what a joint law says about each variable from what it says about their dependence, and \omterm{def:qm:robust-statistics-and-heavy-tails:rank}{rank correlations} and tail dependence
are properties of the \omterm{def:qm:robust-statistics-and-heavy-tails:copula}{copula} alone. For the Gaussian \omterm{def:qm:robust-statistics-and-heavy-tails:copula}{copula} with correlation $\rho$, Spearman's rho is $\frac6\pi\arcsin\frac\rho2$, \omterm{def:qm:robust-statistics-and-heavy-tails:rank}{Kendall's
tau} $\frac2\pi\arcsin\rho$, and both \omterm{def:qm:robust-statistics-and-heavy-tails:copula}{tail dependence coefficients} are zero: joint crashes become asymptotically independent, however high
$\rho$ is. Credit portfolio models inherit that property from the Gaussian \omterm{def:qm:robust-statistics-and-heavy-tails:copula}{copula}, a point One Quant Book 6 takes up with the Student-$t$
\omterm{def:qm:robust-statistics-and-heavy-tails:copula}{copula}, whose tail dependence is positive.

The dollar and the pound, both priced in euros, have a daily Pearson correlation of 0.42 over 1999--2026. Spearman's rho is 0.406, exactly
what a Gaussian \omterm{def:qm:robust-statistics-and-heavy-tails:copula}{copula} with that correlation implies (0.406), and \omterm{def:qm:robust-statistics-and-heavy-tails:rank}{Kendall's tau} 0.286 against an implied 0.277. The tails disagree: on the
5\% of days when the dollar fell most against the euro, the pound was in its own worst 5\% on 32\% of them, against 20\% for the Gaussian
\omterm{def:qm:robust-statistics-and-heavy-tails:copula}{copula}; at 1\%, 21\% against 10\% (\cref{fig:qm:robust-statistics-and-heavy-tails:copula}). The average dependence is Gaussian; the joint
extremes are not.

\begin{omfigure}
\begin{tikzpicture}
\begin{axis}[width=7cm,height=7cm,xlabel={rank of USD per EUR return},ylabel={rank of GBP per EUR return},
  tick label style={font=\small},label style={font=\small},xmin=0,xmax=1,ymin=0,ymax=1,
  table/col sep=comma]
\addplot[only marks,mark=*,mark size=0.35pt,omDef!70] table[x=u,y=v]{figdata/methods/15-robust-statistics-and-heavy-tails/copula.csv};
\draw[omThm,thick] (axis cs:0,0) rectangle (axis cs:0.05,0.05);
\draw[omThm,thick] (axis cs:0.95,0.95) rectangle (axis cs:1,1);
\end{axis}
\end{tikzpicture}

\omcaption{The empirical \omterm{def:qm:robust-statistics-and-heavy-tails:copula}{copula} of daily euro returns against the dollar and the pound, 1999--2026 (every fourth day shown): each point
is a day's pair of ranks scaled to $(0, 1)$. The squares mark the joint 5\% tails, where 32\% of the dollar's worst days are also among the
pound's worst, against 20\% under a Gaussian \omterm{def:qm:robust-statistics-and-heavy-tails:copula}{copula} with the same correlation. Data: ECB euro reference rates (source: ECB
statistics).}\label{fig:qm:robust-statistics-and-heavy-tails:copula}
\end{omfigure}

\section{Heavy tails}

\begin{definition}[Heavy-tailed distribution, tail index]\label{def:qm:robust-statistics-and-heavy-tails:heavy}
A distribution is \emph{heavy-tailed}\index{heavy-tailed distribution} (in the Pareto sense) if $\P(X > x) = x^{-\alpha}L(x)$ for a slowly
varying $L$ ($L(tx)/L(x) \to 1$ as $x \to \infty$ for every $t > 0$). The exponent $\alpha > 0$ is its \emph{tail index}\index{tail index}.
\end{definition}

\begin{proposition}[Moments beyond the tail index]\label{prop:qm:robust-statistics-and-heavy-tails:moments}
If $X \ge 0$ has \omterm{def:qm:robust-statistics-and-heavy-tails:heavy}{tail index} $\alpha$, then $\E[X^\beta] < \infty$ for $\beta < \alpha$ and $\E[X^\beta] = \infty$ for $\beta > \alpha$.
\end{proposition}

\begin{proof}
$\E[X^\beta] = \int_0^\infty\beta x^{\beta-1}\P(X > x)\,dx$. For large $x$, a slowly varying $L$ satisfies $x^{-\varepsilon} \le L(x) \le
x^\varepsilon$ for any $\varepsilon > 0$ (Potter's bounds), so the integrand lies between $x^{\beta - \alpha - 1 \mp \varepsilon}$: integrable
at infinity when $\beta < \alpha$ (take $\varepsilon$ small) and not when $\beta > \alpha$.
\end{proof}

A Student $t$ with $\nu$ degrees of freedom has \omterm{def:qm:robust-statistics-and-heavy-tails:heavy}{tail index} $\nu$. Mandelbrot's 1963 study of cotton prices proposed stable laws, whose
variance is infinite; the Hill plot below gives EUR/USD's daily losses a \omterm{def:qm:robust-statistics-and-heavy-tails:heavy}{tail index} near four, so a finite variance and an unreliable
fourth moment, in line with Cont's survey of stylised facts, which finds tail indices above two and below five for most return series studied. The practical consequence is that the kurtosis, a fourth moment, is estimated badly or not at all: EUR/USD's sample
kurtosis over 27 years is 6.3, and a few days decide it. A quantile-quantile plot shows the tails directly
(\cref{fig:qm:robust-statistics-and-heavy-tails:qq}).

\begin{omfigure}
\begin{tikzpicture}
\begin{axis}[width=8cm,height=6.2cm,xlabel={normal quantile},ylabel={standardised EUR/USD return},
  tick label style={font=\small},label style={font=\small},xmin=-4.2,xmax=4.2,ymin=-8.5,ymax=8.5,grid=major,grid style={gray!25},
  legend style={font=\footnotesize,at={(0.03,0.97)},anchor=north west,fill=white,draw=none},table/col sep=comma]
\addplot[only marks,mark=*,mark size=0.9pt,omDef] table[x=normal,y=sample]{figdata/methods/15-robust-statistics-and-heavy-tails/qq.csv};
\addplot[black,dashed,thick,domain=-4.2:4.2,samples=2] {x};
\legend{daily returns 1999--2026,normal law}
\end{axis}
\end{tikzpicture}

\omcaption{Quantile-quantile plot of the 7\,098 daily EUR/USD returns, standardised, against the normal law (every observation in the
tails, every fiftieth in the body). The largest moves, $-8.2$ and $+7.2$ standard deviations, sit where a normal law puts odds
below one in a hundred billion. Data: ECB euro reference rates (source: ECB statistics).}\label{fig:qm:robust-statistics-and-heavy-tails:qq}
\end{omfigure}

\section{Extreme-value theory}

\begin{definition}[Extreme-value theory, GEV distribution]\label{def:qm:robust-statistics-and-heavy-tails:gev}
\emph{Extreme-value theory}\index{extreme-value theory} studies the laws of maxima and of exceedances of high thresholds. The
\emph{generalised extreme value distribution}\index{generalised extreme value distribution} with shape $\xi$ has distribution function
$G_\xi(x) = \exp\bigl(-(1 + \xi x)^{-1/\xi}\bigr)$ for $1 + \xi x > 0$ (and $\exp(-e^{-x})$ for $\xi = 0$), up to location and scale.
\end{definition}

\begin{theorem}[Fisher--Tippett--Gnedenko]\label{thm:qm:robust-statistics-and-heavy-tails:ftg}
If, for iid $X_i$ and constants $a_n > 0$, $b_n$, the law of $(\max_{i \le n}X_i - b_n)/a_n$ converges to a non-degenerate limit, that limit is a
\omterm{def:qm:robust-statistics-and-heavy-tails:gev}{generalised extreme value distribution}. Heavy tails with index $\alpha$ give $\xi = 1/\alpha > 0$; normal, exponential and similar tails
give $\xi = 0$; bounded tails give $\xi < 0$.
\end{theorem}

\begin{definition}[Generalised Pareto distribution, peaks over threshold]\label{def:qm:robust-statistics-and-heavy-tails:gpd}
The \emph{generalised Pareto distribution}\index{generalised Pareto distribution} with shape $\xi$ and scale $\beta$ has $\P(Y > y) = (1 +
\xi y/\beta)^{-1/\xi}$ for $y \ge 0$ (and $e^{-y/\beta}$ for $\xi = 0$). The \emph{peaks-over-threshold}\index{peaks over threshold} method fits
it by maximum likelihood to the excesses $X - u$ of the observations above a high threshold $u$.
\end{definition}

\begin{theorem}[Pickands--Balkema--de Haan]\label{thm:qm:robust-statistics-and-heavy-tails:pbdh}
For a distribution in the domain of attraction of $G_\xi$, the law of the excess $X - u$ given $X > u$ is approximated, uniformly in
the excess, by a generalised Pareto law with shape $\xi$ and some scale $\beta(u)$, the error vanishing as $u$ tends to the upper end
of the support.
\end{theorem}

The theorem licenses an extrapolation: with $N_u$ of $n$ observations above $u$, the loss exceeded with probability $1 - p$ is
\[
\mathrm{VaR}_p = u + \frac\beta\xi\Bigl(\bigl(\tfrac{N_u/n}{1 - p}\bigr)^\xi - 1\Bigr).
\]
For EUR/USD daily losses above their 95\% quantile ($u = 0.93\%$, 355 excesses), maximum likelihood gives $\xi = 0.069$ and $\beta = 0.34\%$,
and the one-in-a-thousand-days loss is 2.47\%. The fitted $\xi$ depends on the threshold (0.03 above the 90\% quantile, 0.21 above the
99\%), the quantile much less (2.43\% to 2.49\% over the same thresholds): the extrapolation is steadier than the parameter.

\section{Estimating the tail index}

\begin{definition}[Hill estimator]\label{def:qm:robust-statistics-and-heavy-tails:hill}
With the order statistics $X_{(1)} \ge X_{(2)} \ge \dots$ of a positive sample, the \emph{Hill estimator}\index{Hill estimator} of the \omterm{def:qm:robust-statistics-and-heavy-tails:heavy}{tail index}
from the $k$ largest is $\hat\alpha_k = \bigl(\frac1k\sum_{i=1}^k\ln X_{(i)} - \ln X_{(k+1)}\bigr)^{-1}$ (Hill, 1975).
\end{definition}

It is the \omterm{def:qm:estimation:mle}{maximum likelihood estimator} of $\alpha$ for an exact Pareto tail above $X_{(k+1)}$, and when $k \to \infty$ with $k/n \to 0$ at a
suitable rate, $\sqrt k(\hat\alpha_k - \alpha) \to \mathcal N(0, \alpha^2)$. The choice of $k$ is a bias-variance trade: small $k$ is noisy,
large $k$ reaches into the body of the law, where the Pareto form fails. The Hill plot shows the trade
(\cref{fig:qm:robust-statistics-and-heavy-tails:hill}): for EUR/USD losses, $\hat\alpha$ is 4.9 at $k = 50$ ($\pm 0.7$), near 3.9 from $k =
100$ to 200, and drifts down to 2.6 at $k = 700$, where the body is being fitted.

\begin{omfigure}
\begin{tikzpicture}
\begin{axis}[width=10cm,height=5cm,xlabel={number of largest losses used, $k$},ylabel={Hill $\hat\alpha_k$},
  tick label style={font=\small},label style={font=\small},xmin=20,xmax=700,ymin=0,ymax=8,grid=major,grid style={gray!25},
  legend style={font=\footnotesize,at={(0.97,0.97)},anchor=north east,fill=white,draw=none},table/col sep=comma]
\addplot[name path=hi,draw=none,forget plot] table[x=k,y=hi]{figdata/methods/15-robust-statistics-and-heavy-tails/hill.csv};
\addplot[name path=lo,draw=none,forget plot] table[x=k,y=lo]{figdata/methods/15-robust-statistics-and-heavy-tails/hill.csv};
\addplot[omDef!20,forget plot] fill between[of=hi and lo];
\addplot[omDef,very thick] table[x=k,y=alpha]{figdata/methods/15-robust-statistics-and-heavy-tails/hill.csv};
\legend{Hill estimate ($\pm 2$ se shaded)}
\end{axis}
\end{tikzpicture}

\omcaption{Hill plot of the daily losses of EUR/USD, 1999--2026: the \omterm{def:qm:robust-statistics-and-heavy-tails:heavy}{tail index} estimated from the $k$ largest losses, with two asymptotic
\omterm{def:qm:estimation:estimator}{standard errors}. A plateau near 3.9 for $k$ between 100 and 200; beyond, the estimate drifts as the body of the distribution enters.
Data: ECB euro reference rates (source: ECB statistics).}\label{fig:qm:robust-statistics-and-heavy-tails:hill}
\end{omfigure}

Three models then answer the risk manager's question, the one-in-a-thousand-days loss (\cref{fig:qm:robust-statistics-and-heavy-tails:tail}).
The normal law with the sample volatility of 0.58\% puts it at 1.80\%, and 35 days of the 7\,098 lost more, against 7 expected. A Student $t$
fitted by maximum likelihood ($\nu = 4.7$, scale 0.44\%) puts it at 2.74\%. The generalised Pareto fit puts it at 2.47\%, exceeded on 6 days;
the empirical quantile is 2.27\%.

\begin{omfigure}
\begin{tikzpicture}
\begin{semilogyaxis}[width=10cm,height=5.4cm,xlabel={daily loss of EUR/USD (\%)},ylabel={probability of a larger loss},
  tick label style={font=\small},label style={font=\small},xmin=0.5,xmax=4.5,ymin=1e-5,ymax=0.5,grid=major,grid style={gray!25},
  unbounded coords=jump,
  legend columns=4,legend style={font=\footnotesize,at={(0.5,-0.3)},anchor=north,draw=none,fill=none,
  /tikz/every even column/.append style={column sep=6pt}},table/col sep=comma]
\addplot[only marks,mark=*,mark size=1.3pt,black] table[x=loss,y=empirical]{figdata/methods/15-robust-statistics-and-heavy-tails/tail.csv};
\addplot[omThm,very thick,dashed] table[x=loss,y=normal]{figdata/methods/15-robust-statistics-and-heavy-tails/tail.csv};
\addplot[black!60,very thick,dotted] table[x=loss,y=student]{figdata/methods/15-robust-statistics-and-heavy-tails/tail.csv};
\addplot[omDef,very thick] table[x=loss,y=gpd]{figdata/methods/15-robust-statistics-and-heavy-tails/tail.csv};
\addplot[black!40,thin,sharp plot] coordinates {(0.5,0.001) (4.5,0.001)};
\legend{data,normal,Student $t$,{GPD above 0.93\%}}
\end{semilogyaxis}
\end{tikzpicture}

\omcaption{Exceedance probabilities of daily EUR/USD losses, 1999--2026, against three fits: the normal law, a Student $t$ ($\nu = 4.7$) and a
generalised Pareto law above the 95\% quantile. The grey line is the one-in-a-thousand level, crossed at 1.80\%, 2.74\% and 2.47\%. Data:
ECB euro reference rates (source: ECB statistics).}\label{fig:qm:robust-statistics-and-heavy-tails:tail}
\end{omfigure}

Now put the bad print of the hook into the full history. The normal quantile rises to 2.27\% (+26\%), because one pair of $\pm 27\%$ returns
inflates the standard deviation by a quarter. The generalised Pareto quantile rises to 2.85\% (+15\%): the bad loss is the largest excess by
far and lifts the fitted shape from 0.07 to 0.22. The Student $t$ moves least, to 2.91\% (+6\%), its likelihood bounding the influence of any
single point. No tail estimate is robust in the sense of the first section; the tail is where the information is scarce, and one
observation is a large share of it. The defence is upstream, in the filter.

\section{Tutorial: the bad tick}\label{tut:qm:robust-statistics-and-heavy-tails:badtick}

\begin{tutorial}
\textbf{Goal.} Estimate EUR/USD's volatility and tail robustly, then corrupt one print and measure what moves.
\textbf{End state:} \cref{fig:qm:robust-statistics-and-heavy-tails:sensitivity,fig:qm:robust-statistics-and-heavy-tails:hill,fig:qm:robust-statistics-and-heavy-tails:tail}
and the three one-in-a-thousand-days losses with and without the bad print.
\begin{tutsteps}
\item \textbf{The \omterm{def:qm:robust-statistics-and-heavy-tails:hill}{Hill estimator}} from the $k$ largest losses.
\omcode{code/firm/robust/firm_robust.py}{124}{132}{The Hill estimator with its asymptotic standard error.}\label{lst:qm:robust-statistics-and-heavy-tails:hill}
\item \textbf{The generalised Pareto likelihood, fit and quantile.}
\omcode{code/firm/robust/firm_robust.py}{166}{192}{Peaks over threshold: likelihood, fit and quantile.}\label{lst:qm:robust-statistics-and-heavy-tails:gpd}
\item \textbf{Run} \texttt{scale\_table(False)}, \texttt{scale\_table(True)}, \texttt{tail\_quantiles(False)}, \texttt{tail\_quantiles(True)} and
\texttt{dependence()} in \texttt{qm\_robust.py}, then \texttt{fig\_robust.py}.
\end{tutsteps}
\textbf{What to change next.} Build the filter: flag a print whose return exceeds ten MADs of the trailing 60 days and is reversed the next
day, and check that it catches the bad print without flagging 2008; estimate the upper tail (gains) and compare.
\end{tutorial}

\section{Build: the robust statistics module}\label{bld:qm:robust-statistics-and-heavy-tails:robust}

\begin{build}
\bfield{Purpose} Robust scales for every volatility the miniature firm estimates from prints, and the tail estimates its risk limits rest on.
\bfield{Interface} \texttt{mad}, \texttt{qn\_scale}, \texttt{huber\_location}; \texttt{winsorize}, \texttt{trimmed\_mean}; \texttt{spearman},
\texttt{kendall\_tau}, \texttt{tail\_dependence(x, y, q)}; \texttt{hill(x, k)}; \texttt{gpd\_fit(excess)}, \texttt{gpd\_quantile(u, xi, beta,
p\_u, p)}.
\bfield{Rules} Scales are reported with the \omterm{def:qm:estimation:estimator}{estimator} named; a tail quantile is reported with its threshold, its number of excesses and its
sensitivity to the threshold; no $O(n^2)$ memory for pairwise statistics.
\bfield{Acceptance tests} \texttt{code/firm/robust/tests/}: MAD and $Q_n$ consistent at the normal and stable under 10\% gross errors; $Q_n$ equal
to its brute-force definition; the rank-correlation formulas of the Gaussian \omterm{def:qm:robust-statistics-and-heavy-tails:copula}{copula}; Hill on an exact Pareto; the GPD fit recovers its
parameters and its quantile is exceeded at the stated rate.
\bfield{Stretch} Block maxima with the GEV; a declustered \omterm{def:qm:robust-statistics-and-heavy-tails:gpd}{peaks-over-threshold} for dependent losses; the Student-$t$ \omterm{def:qm:robust-statistics-and-heavy-tails:copula}{copula} fitted to the
dollar and the pound.
\end{build}

\begin{omsources}
\item R.~Cont, ``Empirical properties of asset returns: stylized facts and statistical issues'', \emph{Quantitative Finance} 1, 2001.
\item P.~J.~Huber, ``Robust estimation of a location parameter'', \emph{Annals of Mathematical Statistics} 35, 1964.
\item F.~R.~Hampel, ``The influence curve and its role in robust estimation'', \emph{Journal of the American Statistical Association} 69, 1974.
\item P.~J.~Rousseeuw and C.~Croux, ``Alternatives to the median absolute deviation'', \emph{Journal of the American Statistical Association}
88, 1993.
\item A.~Sklar, ``Fonctions de répartition à $n$ dimensions et leurs marges'', \emph{Publications de l'Institut de Statistique de
l'Université de Paris} 8, 1959.
\item R.~A.~Fisher and L.~H.~C.~Tippett, ``Limiting forms of the frequency distribution of the largest or smallest member of a sample'',
\emph{Proceedings of the Cambridge Philosophical Society} 24, 1928; B.~Gnedenko, \emph{Annals of Mathematics} 44, 1943.
\item A.~A.~Balkema and L.~de Haan, ``Residual life time at great age'', \emph{Annals of Probability} 2, 1974; J.~Pickands, ``Statistical
inference using extreme order statistics'', \emph{Annals of Statistics} 3, 1975.
\item B.~M.~Hill, ``A simple general approach to inference about the tail of a distribution'', \emph{Annals of Statistics} 3, 1975.
\item B.~Mandelbrot, ``The variation of certain speculative prices'', \emph{Journal of Business} 36, 1963.
\item European Central Bank, euro foreign exchange reference rates against the dollar and the pound (ECB Data Portal, series
\texttt{EXR.D.USD.EUR.SP00.A} and \texttt{EXR.D.GBP.EUR.SP00.A}), accessed 24 September 2026.
\end{omsources}

\section{Exercises}

\begin{exercise}[$\star$]\label{exo:qm:robust-statistics-and-heavy-tails:1}
What are the \omterm{def:qm:robust-statistics-and-heavy-tails:influence}{breakdown points} of the 10\%-trimmed mean and of the interquartile range?
\end{exercise}

\begin{exercise}[$\star$]\label{exo:qm:robust-statistics-and-heavy-tails:2}
A Pareto tail has index 3. Which of the mean, the variance and the kurtosis exist?
\end{exercise}

\begin{exercise}[$\star$]\label{exo:qm:robust-statistics-and-heavy-tails:3}
For a Gaussian \omterm{def:qm:robust-statistics-and-heavy-tails:copula}{copula} with correlation 0.42, compute Spearman's rho and \omterm{def:qm:robust-statistics-and-heavy-tails:rank}{Kendall's tau}.
\end{exercise}

\begin{exercise}[$\star\star$]\label{exo:qm:robust-statistics-and-heavy-tails:4}
Show that the Huber estimating function with $c = 1.345$ has 95\% efficiency at the normal law: evaluate $(2\Phi(c) - 1)^2/\E[\psi_c(Z)^2]$.
\end{exercise}

\begin{exercise}[$\star\star$]\label{exo:qm:robust-statistics-and-heavy-tails:5}
Show that the \omterm{def:qm:robust-statistics-and-heavy-tails:hill}{Hill estimator} is the \omterm{def:qm:estimation:mle}{maximum likelihood estimator} of $\alpha$ for the $k$ largest observations of an exact Pareto law above
$X_{(k+1)}$.
\end{exercise}

\begin{exercise}[$\star\star$]\label{exo:qm:robust-statistics-and-heavy-tails:6}
With $u = 0.93\%$, 355 excesses out of 7\,098 days, $\xi = 0.069$ and $\beta = 0.34\%$, compute the loss exceeded once in ten thousand days.
\end{exercise}

\begin{exercise}[$\star\star\star$]\label{exo:qm:robust-statistics-and-heavy-tails:7}
\emph{Coding.} Refit the generalised Pareto law to EUR/USD losses above the 90\%, 95\%, 97.5\% and 99\% quantiles, and tabulate $\xi$ and the
one-in-a-thousand-days loss.
\end{exercise}

\begin{exercise}[$\star\star\star$]\label{exo:qm:robust-statistics-and-heavy-tails:8}
\emph{Find the flaw.} ``Our risk model estimates each currency's volatility with the standard deviation of the last 250 days, but we
winsorise the returns at the 1\% and 99\% quantiles first, so bad prints cannot hurt us.''
\end{exercise}

\section{Problem: The Bad Tick}

\begin{problem}[{Weekend problem --- one erroneous print and the one-in-a-thousand-days loss}]\label{pb:qm:robust-statistics-and-heavy-tails:1}
The data are the ECB's daily euro reference rates against the dollar, 4 January 1999 to 23 September 2026 (7\,098 daily returns). In the
corrupted version, the print of 24 March 2026 reads 1.5172 instead of 1.1572.

\medskip\noindent\textbf{Part I --- Scale.}
\begin{enumerate}
\item What log-price error does the bad print make, and which returns does it affect?
\item What are the standard deviation, MAD and $Q_n$ of the last 250 returns, clean and corrupted?
\item What are the mean, the median and the Huber location of those returns, clean and corrupted?
\item How many days in the full history lost more than four standard deviations, against a normal law's expectation?
\item What is the sample kurtosis, and why is it a poor summary?
\end{enumerate}

\medskip\noindent\textbf{Part II --- The tail.}
\begin{enumerate}[resume]
\item What does the Hill plot give for the \omterm{def:qm:robust-statistics-and-heavy-tails:heavy}{tail index}, and where?
\item What are the threshold, the number of excesses and the fitted $\xi$ and $\beta$ of the \omterm{def:qm:robust-statistics-and-heavy-tails:gpd}{peaks-over-threshold} fit?
\item What degrees of freedom and scale does the Student-$t$ fit give?
\item What are the three one-in-a-thousand-days losses, and the empirical one?
\item How many days exceeded the normal and the GPD quantiles, against the 7 expected?
\end{enumerate}

\medskip\noindent\textbf{Part III --- The bad print in the tail.}
\begin{enumerate}[resume]
\item What do the three quantiles become with the bad print?
\item Why does the normal quantile move most?
\item What happens to the fitted $\xi$?
\item Why does the Student $t$ move least?
\item Would a robust scale in the normal model have saved it?
\end{enumerate}

\medskip\noindent\textbf{Part IV --- Judgement.}
\begin{enumerate}[resume]
\item Which quantile would you put in the risk system, and with what caveat?
\item Where should the bad print have been stopped?
\item How does the dollar's tail dependence with the pound change a two-currency limit?
\item State the \emph{named result}: the three quantiles and the error one bad print introduces into each.
\item In one sentence: what can and cannot be made robust?
\end{enumerate}
\end{problem}

\section{Interview questions}

\begin{interviewq}[$\star$ \iqroles{researcher, developer}]\label{iq:qm:robust-statistics-and-heavy-tails:1}
Your volatility estimate for a quiet currency jumps sevenfold overnight. What do you check first?
\end{interviewq}

\begin{interviewq}[$\star\star$ \iqroles{researcher, risk}]\label{iq:qm:robust-statistics-and-heavy-tails:2}
What is a \omterm{def:qm:robust-statistics-and-heavy-tails:influence}{breakdown point}? Give the \omterm{def:qm:robust-statistics-and-heavy-tails:influence}{breakdown points} of the mean, the median and the MAD.
\end{interviewq}

\begin{interviewq}[$\star\star$ \iqroles{researcher, risk}]\label{iq:qm:robust-statistics-and-heavy-tails:3}
What does a \omterm{def:qm:robust-statistics-and-heavy-tails:heavy}{tail index} of 3 imply for the moments of returns, and for estimating kurtosis?
\end{interviewq}

\begin{interviewq}[$\star\star$ \iqroles{risk, researcher}]\label{iq:qm:robust-statistics-and-heavy-tails:4}
How would you estimate a 99.9\% daily loss quantile from 25 years of data?
\end{interviewq}

\begin{interviewq}[$\star\star$ \iqroles{researcher, mle}]\label{iq:qm:robust-statistics-and-heavy-tails:5}
Why use \omterm{def:qm:robust-statistics-and-heavy-tails:rank}{rank correlation} instead of Pearson correlation, and what does neither tell you?
\end{interviewq}

\begin{interviewq}[$\star\star\star$ \iqroles{researcher, risk}]\label{iq:qm:robust-statistics-and-heavy-tails:6}
What does Sklar's theorem say, and why does the Gaussian \omterm{def:qm:robust-statistics-and-heavy-tails:copula}{copula} understate joint crashes?
\end{interviewq}
